When a convolutional network is trained end to end through a variational quantum circuit
(VQC) classifier head, the head shapes the features the network learns. We measure this
effect with one common linear probe on the frozen 8-dimensional bottleneck of every trained
network. On CIFAR-10 with 4,000 training images, an 8-qubit VQC head with CNOT rings leaves
a bottleneck whose probe accuracy is 4.7 points below that of the same network trained with
a linear head (5 seeds, lower for every seed); removing the CNOT rings recovers 3.8 of those
points. An exact identity organizes the pattern: with every variational rotation at zero,
each readout is a product of cosines of the encoded inputs over a set fixed by the CNOT
layers. Across eleven circuit and classical heads, probe accuracy falls as the mean order of
these products rises, from 73.6 at order 1 to 69.8 at the default order 2.5 and 60.1 at
order 5.25. Classical heads trained on the same products, with no circuit, follow the same
trend, and a circuit that keeps its CNOTs but has product order 1 recovers 3.1 points.
Without CNOTs the head is exactly a classical per-coordinate sinusoidal model. Under a
revised protocol with longer training, validation-based selection and per-head
learning-rate tuning, the VQC's apparent advantage over a small MLP on MNIST shrinks from up
to 18 points to at most 3.3, and among linear, MLP and default VQC heads the linear head has
the highest mean test accuracy in all six settings.
